\chapter{Fractions, racines et puissances}

% ══════════════════════════════════════════════
\section{Fractions}
% ══════════════════════════════════════════════

\subsection{Calculs avec les fractions}

Y'a rien qui va

% \begin{prop}[Opérations sur les fractions]
% 	\bi
% 	\item Somme ou différence
% 	$$\frac{a}{D} + \frac{b}{D} = \frac{a+b}{D}\qquad\qquad\frac{a}{D} - \frac{b}{D} = \frac{a-b}{D}$$
% 	\item Produit ou quotient
% 	$$\frac{a}{b} \times \frac{c}{d} = \frac{a \times c}{b \times d}\qquad\qquad\frac{a}{b} \div \frac{c}{d} = \frac{a}{b} \times \frac{d}{c}$$
% 	\ei
% \end{prop}
%
% % \begin{methode}[Effectuer des calculs de fractions]
% % 	Calculer :
% %
% % 	$$A = \dfrac{5}{4} + \dfrac{6}{16} \qquad B = \dfrac{5}{3} - \dfrac{6}{5} \qquad C = \dfrac{2}{-3} \times \dfrac{-5}{11}$$
% %
% % 	$$D = \dfrac{3}{4} \div \dfrac{-5}{8} \qquad E = \dfrac{8}{7} - \dfrac{4}{7} \times \dfrac{5}{3}$$
% %
% % 	\begin{correction}
% % 		$\begin{array}{r|r|r|r|r}
% % 				\def\arraystretch{2}
% % 				\begin{array}{rcl}
% % 					A & = & \dfrac{5}{4} + \dfrac{6}{16}   \\
% % 					  & = & \dfrac{20}{16} + \dfrac{6}{16} \\
% % 					  & = & \dfrac{26}{16}                 \\
% % 					  & = & \dfrac{13}{8}
% % 				\end{array}
% % 				 &
% % 				\def\arraystretch{2}
% % 				\begin{array}{rcl}
% % 					B & = & \dfrac{5}{3} - \dfrac{6}{5}     \\
% % 					  & = & \dfrac{25}{15} - \dfrac{18}{15} \\
% % 					  & = & \dfrac{7}{15}                   \\&&
% % 				\end{array}
% % 				 &
% % 				\def\arraystretch{2}
% % 				\begin{array}{rcl}
% % 					C & = & \dfrac{2}{-3} \times \dfrac{-5}{11}   \\
% % 					  & = & \dfrac{2 \times (-5)}{(-3) \times 11} \\
% % 					  & = & \dfrac{-10}{-33}                      \\
% % 					  & = & \dfrac{10}{33}
% % 				\end{array}
% % 				 &
% % 				\def\arraystretch{2}
% % 				\begin{array}{rcl}
% % 					D & = & \dfrac{3}{4} \div \dfrac{-5}{8}   \\
% % 					  & = & \dfrac{3}{4} \times \dfrac{8}{-5} \\
% % 					  & = & \dfrac{24}{-20}                   \\
% % 					  & = & -\dfrac{6}{5}
% % 				\end{array}
% % 				 &
% % 				\def\arraystretch{2}
% % 				\begin{array}{rcl}
% % 					E & = & \dfrac{8}{7} - \dfrac{4}{7} \times \dfrac{5}{3} \\
% % 					  & = & \dfrac{8}{7} - \dfrac{20}{21}                   \\
% % 					  & = & \dfrac{24}{21} - \dfrac{20}{21}                 \\
% % 					  & = & \dfrac{4}{21}
% % 				\end{array}
% % 			\end{array}$
% % 	\end{correction}
% % \end{methode}
%
% \subsection{Réduire au même dénominateur}
%
% \begin{prop}[Réduction au même dénominateur]
% 	\[
% 		\frac{a}{b} + \frac{c}{d} = \frac{ad}{bd} + \frac{cb}{bd} = \frac{ad + cb}{bd}
% 	\]
% \end{prop}
% % \newpage
% % \begin{methode}[Réduire des expressions au même dénominateur]
% % 	Réduire au même dénominateur :
% % 	\[
% % 		A = \frac{7}{x-2} - \frac{5}{3}
% % 		\qquad\qquad
% % 		B = 3 + \frac{5x}{2x+1}
% % 	\]
% % 	\begin{correction}
% % 		$\begin{array}{c|c}
% % 				\def\arraystretch{2}
% % 				\begin{array}{rcl}
% % 					A & = & \dfrac{7}{x-2} - \dfrac{5}{3}                                 \\
% % 					  & = & \dfrac{7 \times 3}{3(x-2)} - \dfrac{5(x-2)}{3(x-2)}           \\
% % 					  & = & \dfrac{21 - 5(x-2)}{3(x-2)}                                   \\
% % 					  & = & \dfrac{21 - 5x + 10}{3(x-2)} \qquad = \dfrac{31 - 5x}{3(x-2)} \\
% % 				\end{array}\qquad & \qquad
% % 				\def\arraystretch{2}
% % 				\begin{array}{rcl}
% % 					B & = & 3 + \dfrac{5x}{2x+1}                     \\
% % 					  & = & \dfrac{3(2x+1)}{2x+1} + \dfrac{5x}{2x+1} \\
% % 					  & = & \dfrac{6x + 3 + 5x}{2x+1}                \\
% % 					  & = & \dfrac{11x + 3}{2x+1}                    \\
% % 				\end{array}
% % 			\end{array}$
% % 	\end{correction}
% % \end{methode}
% %
% % ══════════════════════════════════════════════
% \section{Puissances}
% % ══════════════════════════════════════════════
%
% \subsection{Rappels}
%
% \begin{defin}{Puissance entière}
% 	Pour tout nombre $a$ non nul et tout entier $n \ge 1$ :
% 	\[
% 		a^n = \underbrace{a \times a \times \cdots \times a}_{n \text{ fois}}
% 	\]
% 	Par convention : $a^1 = a$ \quad $a^0 = 1$ \quad $0^n = 0$ \quad $1^n = 1$
% \end{defin}
%
% \begin{rem}
% 	Ne pas confondre :
% 	\[
% 		(-3)^3 = (-3)\times(-3)\times(-3) = 27
% 		\quad \text{et} \quad
% 		-3^3 = -(3\times 3 \times 3) = -27
% 	\]
% \end{rem}
%
% % \begin{exercice}[Règle des signes]
% % 	Calculer : $(-5)^2$ ; $-1^2$ ; $(-1)^2$ ; $-3^3$ ; $(-2)^2$ ; $-7^2$ ; $(-9)^0$ ; $-9^0$
% %
% % 	\begin{correction}
% % 		\begin{itemize}
% % 			\begin{multicols}{2}
% % 				\item $(-5)^2=25$
% % 				\item $-1^2=-1$
% % 				\item $(-1)^2=1$
% % 				\item $-3^3=-9$
% % 				\item $(-2)^2=4$
% % 				\item $-7^2=-49$
% % 				\item $(-9)^0=1$
% % 				\item $-9^0=-1$
% % 			\end{multicols}
% % 		\end{itemize}
% % 	\end{correction}
% % \end{exercice}
% % \newpage
% \subsection{Opérations sur les puissances}
%
% \begin{prop}[Règles de calcul ($n, p\in\Z$ et $a \neq 0$, $b \neq 0$)]
% 	% \begin{itemize}[itemsep=3mm]
% 	\begin{itemize}
% 		\item \textbf{Multiplication :}\qquad $a^n \times a^p = a^{n+p}$
% 		\item \textbf{Division :}\qquad $\dfrac{a^n}{a^p} = a^{n-p}$
% 		\item \textbf{Inverses :}\qquad $a^{-1} = \dfrac{1}{a} \quad \text{et} \quad a^{-n} = \dfrac{1}{a^n}$
% 		\item \textbf{Puissance de puissance :}\qquad $\left(a^n\right)^p = a^{n \times p}$
% 		\item \textbf{Distribution :}\qquad $(a \times b)^n = a^n \times b^n$
% 	\end{itemize}
% \end{prop}
% % \begin{methode}[Exprimer sous la forme d'une seule puissance]
% % 	Exprimer sous la forme d'une seule puissance :
% % 	\[
% % 		A = \frac{1}{4^2}
% % 		\qquad
% % 		B = 4^5 \times 4^7
% % 		\qquad
% % 		C = \frac{5^4}{5^6}
% % 		\qquad
% % 		D = 7^3 \times \left(7^2\right)^6
% % 		\qquad
% % 		E = 6^7 \times 9^7
% % 	\]
% % 	\begin{correction}
% % 		\begin{itemize}[itemsep=2mm]
% % 			\item $A = \dfrac{1}{4^2} = 4^{-2}$
% % 			\item $B = 4^5 \times 4^7 = 4^{5+7} = 4^{12}$
% % 			\item $C = \dfrac{5^4}{5^6} = 5^{4-6} = 5^{-2}$
% % 			\item $D = 7^3 \times \left(7^2\right)^6 = 7^3 \times 7^{12} = 7^{3+12} = 7^{15}$
% % 			\item $E = 6^7 \times 9^7 = (6 \times 9)^7 = 54^7$
% % 		\end{itemize}
% % 	\end{correction}
% % \end{methode}
% %
% % \begin{methode}[Puissances de 10 et notation scientifique]
% % 	\textbf{a)} Écrire sous la forme $10^n$ :
% % 	\[
% % 		A = 10^4 \times 10^7
% % 		\qquad
% % 		B = \frac{10^{-4}}{10^5}
% % 		\qquad
% % 		C = \left(10^2\right)^{-6}
% % 		\qquad
% % 		D = 10^{-4} \times \left(10^3\right)^{-1}
% % 	\]
% %
% % 	\textbf{b)} Écrire en notation scientifique :
% % 	\bi[itemsep=2mm]
% % 	\begin{multicols}{2}
% % 		\item $A = 4 \times 7 \times 10^{-5} \times 10^{-8}$
% % 		\item $B = \dfrac{7 \times 10^{-4} \times 5 \times 10^8}{56 \times 10^{-9}}$
% % 		\item $C = \dfrac{32 \times 10^{-4} + 6 \times 10^{-3}}{2 \times 10^{-5}}$
% % 	\end{multicols}
% % 	\ei
% %
% % 	\begin{correction}
% %
% % 		\textbf{a)}
% % 		\bi
% % 		\begin{multicols}{2}
% % 			\item $A = 10^{4+7} = 10^{11}$
% % 			\item $B = 10^{-4-5} = 10^{-9}$
% % 			\item $C = 10^{2\times(-6)} = 10^{-12}$
% % 			\item $D = 10^{-4} \times 10^{-3} = 10^{-7}$
% % 		\end{multicols}
% % 		\ei
% % 		\newpage
% % 		\textbf{b)}
% %
% % 		$\begin{array}{r|r|r}
% % 				\def\arraystretch{1.9}
% % 				\begin{array}{rcl}
% % 					A & = & 28 \times 10^{-13}    \\
% % 					  & = & 2{,}8 \times 10^{-12} \\
% % 					  &   &                       \\
% % 					  &   &                       \\
% % 				\end{array}                                                                                            &
% % 				\def\arraystretch{1.9}
% % 				\begin{array}{rcl}
% % 					B & = & \dfrac{7 \times 5}{56} \times \dfrac{10^{-4} \times 10^8}{10^{-9}} \\
% % 					  & = & 0{,}625 \times 10^{-4+8+9}                                         \\
% % 					  & = & 0{,}625 \times 10^{13}                                             \\
% % 					  & = & 6{,}25 \times 10^{12}                                              \\
% % 				\end{array} &
% % 				\def\arraystretch{1.9}
% % 				\begin{array}{rcl}
% % 					C & = & \dfrac{0{,}0032 + 0{,}006}{2 \times 10^{-5}} \\
% % 					  & = & \dfrac{0{,}0092}{2 \times 10^{-5}}           \\
% % 					  & = & 0{,}0046 \times 10^5                         \\
% % 					  & = & 4{,}6 \times 10^{2}                          \\
% % 				\end{array}
% % 			\end{array}$
% % 	\end{correction}
% % \end{methode}
%
% % ══════════════════════════════════════════════
% \section{Racines carrées}
% % ══════════════════════════════════════════════
%
% \subsection{Définition}
%
% \begin{defin}{Racine carrée}
% 	La \textbf{racine carrée} de $a$ ($a \ge 0$) est l'unique nombre positif dont le carré est $a$. On la note $\sqrt{a}$.
% \end{defin}
%
% % \begin{ex}[Racines carrées]
% % 	\bi
% % 	\begin{multicols}{2}
% % 		\item $3^2 = 9$ donc $\sqrt{9} = 3$ \qquad
% % 		\item $2{,}6^2 = 6{,}76$ donc $\sqrt{6{,}76} = 2{,}6$ \qquad
% % 		\item $\sqrt{2} \approx 1{,}4142$ \qquad
% % 		\item $\sqrt{3} \approx 1{,}732$
% % 	\end{multicols}
% % 	\ei
% % \end{ex}
% \begin{prop}[Carrés parfaits]
% 	\[
% 		\rnc\arraystretch{1.5}
% 		\begin{array}{l|l|l}
% 			\sqrt{0}=0  & \sqrt{25}=5 & \sqrt{100}=10 \\
% 			\sqrt{1}=1  & \sqrt{36}=6 & \sqrt{121}=11 \\
% 			\sqrt{4}=2  & \sqrt{49}=7 & \sqrt{144}=12 \\
% 			\sqrt{9}=3  & \sqrt{64}=8 & \sqrt{169}=13 \\
% 			\sqrt{16}=4 & \sqrt{81}=9 &               \\
% 		\end{array}
% 	\]
% \end{prop}
%
% \begin{rem}
% 	La racine carrée d'un nombre \textbf{négatif} n'existe pas : un carré est toujours positif, donc $\sqrt{-5}$ n'existe pas.
% \end{rem}
% % \newpage
% \subsection{Propriétés}
%
% \begin{prop}[Règles de calcul ($a \ge 0$, $b \ge 0$, $b \neq 0$)]
% 	\begin{itemize}[itemsep=3mm]
% 		\item \textbf{Produit :}\qquad $\sqrt{a} \times \sqrt{b} = \sqrt{a \times b}$
% 		\item \textbf{Quotient :}\qquad $\dfrac{\sqrt{a}}{\sqrt{b}} = \sqrt{\dfrac{a}{b}}$
% 		\item \textbf{Carré d'une racine :}\qquad $\left(\sqrt{a}\right)^2 = a$
% 		\item \textbf{Racine d'un carré :}\qquad $\sqrt{a^2} = a$
% 		\item \textbf{Cas particulier :}\qquad $\sqrt{0} = 0 \quad \text{et} \quad \sqrt{1} = 1$
% 	\end{itemize}
% \end{prop}
%
% \begin{attention} De façon générale, $\sqrt{a} + \sqrt{b} \neq \sqrt{a+b}$ et $\sqrt{a} - \sqrt{b} \neq \sqrt{a-b}$.\end{attention}
%
% \begin{demo}[$\sqrt{a} \times \sqrt{b} = \sqrt{a \times b}$]
% 	On a :
% 	\be
% 	\item $\left(\sqrt{a} \times \sqrt{b}\right)^2 = \left(\sqrt{a}\right)^2 \times \left(\sqrt{b}\right)^2 = a \times b$
% 	\item $\left(\sqrt{a \times b}\right)^2 = a \times b$ \quad (car $a$ et $b$ sont positifs)
% 	\ee
% 	Donc $\left(\sqrt{a} \times \sqrt{b}\right)^2 = \left(\sqrt{a \times b}\right)^2$, et donc $\sqrt{a} \times \sqrt{b} = \sqrt{a \times b}$.
% \end{demo}
%
% \begin{demo}[$\sqrt{a} + \sqrt{b} > \sqrt{a+b}$]
% 	\be
% 	\item $\left(\sqrt{a} + \sqrt{b}\right)^2 = \left(\sqrt{a}\right)^2 + 2\sqrt{a}\sqrt{b} + \left(\sqrt{b}\right)^2 = a + b + 2\sqrt{ab}$
% 	\item $\left(\sqrt{a+b}\right)^2 = a + b$
% 	\ee
% 	Donc $\left(\sqrt{a}+\sqrt{b}\right)^2 > \left(\sqrt{a+b}\right)^2$ car $2\sqrt{ab} > 0$, et donc : $$\sqrt{a} + \sqrt{b} > \sqrt{a+b}$$
% \end{demo}
%
% % \begin{methode}[Effectuer des calculs sur les racines carrées]
% % 	Écrire le plus simplement possible :
% % 	\[
% % 		\def\arraystretch{2}
% % 		\begin{array}{lllllll}
% % 			A = \sqrt{32} \times \sqrt{2} \quad & \quad D = \dfrac{\sqrt{98}}{\sqrt{2}} \quad & \quad F = \left(4\sqrt{5}\right)^2                      \\
% % 			B = \sqrt{3} \times \sqrt{27} \quad & \quad E = \dfrac{\sqrt{50}}{\sqrt{72}}\quad & \quad G = \dfrac{\sqrt{32} \times \sqrt{10}}{\sqrt{80}} \\
% % 			C = \sqrt{3} \times \sqrt{36} \times \sqrt{3}
% % 		\end{array}
% % 	\]
% % 	\begin{correction}
% % 		$\def\arraystretch{2}
% % 			\begin{array}{rr}
% % 				\begin{array}{rcl}
% % 					A & = & \sqrt{32 \times 2} = \sqrt{64} = 8                                               \\
% % 					B & = & \sqrt{3 \times 27} = \sqrt{81} = 9                                               \\
% % 					C & = & \sqrt{3 \times 3} \times \sqrt{36} = \sqrt{9} \times \sqrt{36} = 3 \times 6 = 18 \\
% % 					D & = & \sqrt{\dfrac{98}{2}} = \sqrt{49} = 7                                             \\
% % 				\end{array}\qquad & \qquad
% % 				\begin{array}{rcl}
% % 					E & = & \sqrt{\dfrac{50}{72}} = \sqrt{\dfrac{25}{36}} = \dfrac{\sqrt{25}}{\sqrt{36}} = \dfrac{5}{6} \\
% % 					F & = & 4^2 \times \left(\sqrt{5}\right)^2 = 16 \times 5 = 80                                       \\
% % 					G & = & \sqrt{\dfrac{32 \times 10}{80}} = \sqrt{4} = 2                                              \\
% % 					  &   &
% % 				\end{array}
% % 			\end{array}$
% % 	\end{correction}
% % \end{methode}
% % \newpage
% % \subsection{Extraire un carré parfait}
% %
% % \begin{methode}[Écrire sous la forme $a\sqrt{b}$]
% % 	Écrire sous la forme $a\sqrt{b}$, avec $a$ et $b\in\Z$ et $b$ le plus petit possible :
% % 	\[
% % 		A = \sqrt{72}
% % 		\qquad
% % 		B = \sqrt{45}
% % 		\qquad
% % 		C = 3\sqrt{125}
% % 	\]
% %
% % 	\begin{correction}
% % 		$\begin{array}{lll}
% % 				\def\arraystretch{1.3}
% % 				\begin{array}{rcl}
% % 					A & = & \sqrt{72} \\ &=& \sqrt{36 \times 2}\\ &=& \sqrt{36} \times \sqrt{2}\\ &=& 6\sqrt{2}\\
% % 				\end{array} &
% % 				\def\arraystretch{1.3}
% % 				\begin{array}{rcl}
% % 					B & = & \sqrt{45} \\ &=& \sqrt{9 \times 5}\\ &=& \sqrt{9} \times \sqrt{5}\\ &=& 3\sqrt{5} \\
% % 				\end{array}  &
% % 				\def\arraystretch{1.3}
% % 				\begin{array}{rcl}
% % 					C & = & 3\sqrt{125} \\ &=& 3\sqrt{25 \times 5}\\ &=& 3\sqrt{25} \times \sqrt{5}\\ &=& 3 \times 5 \times \sqrt{5} = 15\sqrt{5} \\
% % 				\end{array}
% % 			\end{array}$
% % 		\textit{Pour que $b$ soit le plus petit possible, $b$ ne doit pas contenir de carré parfait.}
% % 	\end{correction}
% % \end{methode}
% %
% % \subsection{Simplifier des expressions contenant des racines}
% %
% % \begin{methode}[Réduire des expressions avec racines carrées]
% % 	\textbf{1)} Écrire le plus simplement possible :
% % 	\bi
% % 	\item $A = 4\sqrt{3} - 2\sqrt{3} + 6\sqrt{3}$
% % 	\item $B = 7\sqrt{2} - 3\sqrt{5} + 8\sqrt{2} - \sqrt{5}$
% % 	\item $C = \left(3 - 2\sqrt{3}\right) - \left(4 - 6\sqrt{3}\right)$
% % 	\ei
% %
% % 	\textbf{2)} Écrire sous la forme $a\sqrt{b}$ ($b$ le plus petit possible) :
% % 	\bi
% % 	\item $D = \sqrt{12} + 7\sqrt{3} - \sqrt{27}$
% % 	\item $E = \sqrt{125} - 2\sqrt{20} + 6\sqrt{80}$
% % 	\ei
% %
% % 	\begin{correction}
% %
% % 		\textbf{1)} On regroupe les membres d'une même « famille » de racines carrées.
% %
% % 		$\rnc{\arraystretch}{2.5}
% % 			\begin{array}{ll}
% % 				\rnc{\arraystretch}{1.5}
% % 				\begin{array}{rcl}
% % 					A & = & 4\sqrt{3} - 2\sqrt{3} + 6\sqrt{3} \\
% % 					  & = & (4 - 2 + 6)\sqrt{3}               \\
% % 					  & = & 8\sqrt{3}
% % 				\end{array}                       &
% % 				\rnc{\arraystretch}{1.5}
% % 				\begin{array}{rcl}
% % 					B & = & 7\sqrt{2} - 3\sqrt{5} + 8\sqrt{2} - \sqrt{5} \\
% % 					  & = & (7 + 8)\sqrt{2} + (-3-1)\sqrt{5}             \\
% % 					  & = & 15\sqrt{2} - 4\sqrt{5}
% % 				\end{array}            \\
% % 				\rnc{\arraystretch}{1.5}
% % 				\begin{array}{rcl}
% % 					C & = & \left(3 - 2\sqrt{3}\right) - \left(4 - 6\sqrt{3}\right) \\
% % 					  & = & 3 - 2\sqrt{3} - 4 + 6\sqrt{3}                           \\
% % 					  & = & -1 + 4\sqrt{3}
% % 				\end{array} & ~
% % 			\end{array}$
% %
% % 		\textbf{2)} On extrait les carrés parfaits pour ramener à la même famille.
% %
% % 		$\rnc{\arraystretch}{1.5}
% % 			\begin{array}{rcl}
% % 				D & = & \sqrt{4 \times 3} + 7\sqrt{3} - \sqrt{9 \times 3}   \\
% % 				  & = & 2\sqrt{3} + 7\sqrt{3} - 3\sqrt{3} \quad = 6\sqrt{3}
% % 			\end{array}\qquad
% % 			\begin{array}{rcl}
% % 				E & = & \sqrt{25 \times 5} - 2\sqrt{4 \times 5} + 6\sqrt{16 \times 5} \\
% % 				  & = & 5\sqrt{5} - 4\sqrt{5} + 24\sqrt{5} \quad = 25\sqrt{5}
% % 			\end{array}$
% % 	\end{correction}
% % \end{methode}
% % \newpage
% % \subsection{Développements avec des racines carrées}
% %
% % \begin{methode}[Développer et réduire avec des racines carrées]
% % 	Développer et réduire :
% % 	\bi
% % 	\begin{multicols}{2}
% % 		\item $A = \left(\sqrt{3} - 4\right)^2$
% % 		\item $B = \left(3 + \sqrt{5}\right)^2$
% % 		\item $C = \left(\sqrt{2} - \sqrt{5}\right)\left(\sqrt{2} + \sqrt{5}\right)$
% % 		\item $D = \left(3 + \sqrt{3}\right)\left(1 - \sqrt{2}\right)$
% % 	\end{multicols}
% % 	\ei
% %
% % 	\begin{correction}
% % 		On applique les identités remarquables et la distributivité, en traitant les racines comme des inconnues.
% %
% % 		$\def\arraystretch{1.4}
% % 			\begin{array}{rcl}
% % 				A & = & \left(\sqrt{3}\right)^2 - 2 \times \sqrt{3} \times 4 + 4^2 \\
% % 				  & = & 3 - 8\sqrt{3} + 16                                         \\
% % 				  & = & 19 - 8\sqrt{3}
% % 			\end{array}
% % 			\qquad
% % 			\def\arraystretch{1.4}
% % 			\begin{array}{rcl}
% % 				B & = & 3^2 + 2 \times 3 \times \sqrt{5} + \left(\sqrt{5}\right)^2 \\
% % 				  & = & 9 + 6\sqrt{5} + 5                                          \\
% % 				  & = & 14 + 6\sqrt{5}
% % 			\end{array}$
% %
% % 		\ms
% %
% % 		$\def\arraystretch{1.4}
% % 			\begin{array}{rcl}
% % 				C & = & \left(\sqrt{2}\right)^2 - \left(\sqrt{5}\right)^2 \\
% % 				  & = & 2 - 5 \qquad = -3
% % 			\end{array}
% % 			\qquad
% % 			\def\arraystretch{1.4}
% % 			\begin{array}{rcl}
% % 				D & = & 3 - 3\sqrt{2} + \sqrt{3} - \sqrt{3} \times \sqrt{2} \\
% % 				  & = & 3 - 3\sqrt{2} + \sqrt{3} - \sqrt{6}
% % 			\end{array}$
% % 	\end{correction}
% % \end{methode}
% %
\end{document}

